\documentclass[11pt]{article}
\usepackage{palestra}
\newcommand{\poc}[3]{\begin{tabular}[b]{@{}c@{}}\bebe{#1}{#2}\\[1mm]#3\end{tabular}}
% una línea del cuaderno de Listillón, para el ojo de Argos
\newcommand{\linea}[1]{\hspace*{4mm}$= #1$\hfill\tikz[baseline=-.6ex]\draw[tinta!45,line width=.5pt](0,0)circle(1.7mm);\\[1.2mm]}
\begin{document}
\begin{ficha}{Potencias}{Su propio frasco}{4}

\begin{encargo}
Elevar al cuadrado es multiplicar un número por sí mismo. En la palestra, el
tirador \textbf{se bebe su propio frasco}: en $(-2)^2$, el $-2$ se bebe un
frasco de $\cdot(-2)$. Al cubo, se lo bebe dos veces.
\end{encargo}

\bloque{1}{Su propio frasco}{En el campo}\quad
{\footnotesize ¿qué se bebe, cuántas veces, y con qué número sale?}

\vspace{1.5mm}
{\small
\begin{tabular}{@{}c@{\hspace{14mm}}c@{\hspace{14mm}}c@{}}
\poc{-2}{\traiciondoble}{\ej{a}$(-2)^2 = $\casillita[7mm]} &
\poc{3}{\triplon}{\ej{b}$(+3)^2 = $\casillita[7mm]} &
\poc{-3}{\traiciontriple}{\ej{c}$(-3)^2 = $\casillita[7mm]}\\[2mm]
\poc{2}{\doblon\doblon}{\ej{d}$(+2)^3 = $\casillita[7mm]} &
\poc{-2}{\traiciondoble\traiciondoble}{\ej{e}$(-2)^3 = $\casillita[7mm]} & \\
\end{tabular}}

\vspace{1.5mm}
{\small En \ej{e}, ¿en qué bando está después del primer frasco? \hueco[1.5cm]\quad
¿Y después del segundo? \hueco[1.5cm]}

\vspace{1.5mm}
\begin{listillon}
Una \textbf{potencia de base negativa} es positiva si el exponente es
\textbf{par}, y negativa si es \textbf{impar}. El exponente solo afecta a lo que
tiene justo debajo: en $(-3)^2$ la base es $-3$, y vale $+9$; en $-3^2$ la base
es 3, y el menos va fuera: $-3^2 = -9$.
\end{listillon}

\bloque{2}{¿Con o sin paréntesis?}{Sobre el papel}\quad
{\footnotesize calcula las dos y pon el signo: $<$, $>$ o $=$}

\vspace{1mm}
{\small\renewcommand{\arraystretch}{2.1}
\begin{tabular}{@{}l@{\hspace{14mm}}l@{}}
\ej{a}$(-3)^2 = $ \casillita[7mm] \ \signo\ \ $-3^2 = $ \casillita[7mm] &
\ej{b}$(-2)^3 = $ \casillita[7mm] \ \signo\ \ $-2^3 = $ \casillita[7mm]\\
\ej{c}$(-1)^4 = $ \casillita[7mm] \ \signo\ \ $-1^4 = $ \casillita[7mm] &
\ej{d}$(-2)^2 = $ \casillita[7mm] \ \signo\ \ $2^2 = $ \casillita[7mm]\\
\end{tabular}}

\vspace{.5mm}
{\small ¿En cuáles da lo mismo con paréntesis que sin él? \hueco[2cm]\quad ¿Por qué? \hueco[4.2cm]}

\bloque{3}{Todo junto}{Sobre el papel}\quad
{\footnotesize primero las potencias, luego los frascos, y al final se suma}

\vspace{1mm}
{\small\renewcommand{\arraystretch}{2.1}
\begin{tabular}{@{}l@{\hspace{10mm}}l@{\hspace{10mm}}l@{}}
\ej{a}$5 - (-2)^2 = $ \casillita[7mm] & \ej{b}$(-3)^2 - 2\cdot 4 = $ \casillita[7mm] &
\ej{c}$-2^2 + (-5)\cdot(-1) = $ \casillita[7mm]\\
\ej{d}$4 + (-2)^3 = $ \casillita[7mm] & \ej{e}$(-1)^3 - (-3)^2 = $ \casillita[7mm] &
\ej{f}$10 - 3\cdot(-2)^2 = $ \casillita[7mm]\\
\end{tabular}}

\alto{$-3^2$ no es $+9$: el cuadrado es del 3, no del $-3$. Primero $3^2 = 9$, y
luego el menos: $-9$. Si se quiere que el $-3$ se beba su frasco, hay que
escribir el paréntesis: $(-3)^2 = +9$.}

\reto{¿Cuánto valen $(-1)^2$, $(-1)^3$, $(-1)^4$…? ¿Y $(-1)^{100}$? ¿Y
$(-1)^{101}$? Explica cómo lo sabes sin hacer cien multiplicaciones.}

\comprueba{Cada código abre esa cuenta en Practicar; si fallas, mira el banquillo.}{%
\qr{1c}{j=potencias\&m=7\&c=n-3_e2}\hfill\qr{1e}{j=potencias\&m=7\&c=n-2_e3}\hfill
\qr{3a}{j=potencias\&m=8\&c=n5_r_n-2_e2}\hfill\qr{3b}{j=potencias\&m=8\&c=n-3_e2_r_q2_n4}\hfill
\qr{3c}{j=potencias\&m=8\&c=r_n2_e2_s_n-5_p-1}\hfill\qr{3d}{j=potencias\&m=8\&c=n4_s_n-2_e3}}

\end{ficha}

% ─────────────────────────────── REVERSO ──────────────────────────────────
\begin{dorso}{Potencias}{El ojo de Argos}{4}

\bloque{1}{El ojo de Argos}{Sobre el papel}\quad
{\footnotesize Listillón se ha equivocado en \textbf{un} paso: marca la
\textbf{primera} línea mal y corrígela al lado}

\vspace{2mm}
{\small
\begin{minipage}[t]{.47\linewidth}
\ej{a}$-8 - 3\cdot(-2)$\\[1.2mm]
\linea{-8 - (-6)}
\linea{-8 - 6}
\linea{-14}
\end{minipage}\hfill
\begin{minipage}[t]{.47\linewidth}
\ej{b}$(-2)^3 + 5$\\[1.2mm]
\linea{(-2)\cdot(-2)\cdot(-2) + 5}
\linea{8 + 5}
\linea{13}
\end{minipage}

\vspace{4mm}
\begin{minipage}[t]{.47\linewidth}
\ej{c}$-3^2 + 2\cdot(-1 + 4)$\\[1.2mm]
\linea{-9 + 2\cdot(-1 + 4)}
\linea{-9 + 2\cdot 5}
\linea{-9 + 10}
\linea{1}
\end{minipage}\hfill
\begin{minipage}[t]{.47\linewidth}
\ej{d}$(-12)\dv (-3) - (-2)^2$\\[1.2mm]
\linea{4 - (-2)^2}
\linea{4 - (+4)}
\linea{4 + 4}
\linea{8}
\end{minipage}}

\vspace{2mm}
{\small ¿Cuál de los fallos de Listillón es el que más se repite en clase? \hueco[5.4cm]}

\bloque{2}{Más todo junto}{Sobre el papel}

\vspace{1mm}
{\small\renewcommand{\arraystretch}{2.3}
\begin{tabular}{@{}l@{\hspace{10mm}}l@{\hspace{10mm}}l@{}}
\ej{a}$2\cdot(-3)^2 - 5\cdot(-1) = $ \casillita[7mm] & \ej{b}$(-4 + 1)^2 - 10 = $ \casillita[7mm] &
\ej{c}$-(-2)^2 - (-2)^3 = $ \casillita[7mm]\\
\ej{d}$(-2)^2\cdot(-3) + 1 = $ \casillita[7mm] & \ej{e}$(-18)\dv (-3)^2 = $ \casillita[7mm] &
\ej{f}$(5 - 7)^3 + (-1)^5 = $ \casillita[7mm]\\
\end{tabular}}

\bloque{3}{Inventa}{Sobre el papel}\quad
{\footnotesize una cuenta tuya, para que la resuelva tu compañero}

\vspace{1.5mm}
{\small Escribe una cuenta con una \textbf{potencia}, un \textbf{paréntesis} y un
\textbf{frasco}, en la que la bandera acabe en $-3$.\\[5mm]
\rule{.7\linewidth}{.4pt}\\[5mm]
Pruébala en la palestra, en Practicar, con «tu cuenta». ¿Ha acabado en $-3$? \hueco[1.5cm]}\par\vspace{4mm}

\reto{Listillón dice que $-2^2$ y $(-2)^2$ son iguales, «que el paréntesis es de
adorno». ¿Con qué número sale a tirar cada uno? Convéncele con el campo.}

\comprueba{Los del bloque 1, en el ojo de Argos de la app (con otro fallo); los del reto, jugados.}{%
\qr{1a}{j=potencias\&m=9\&c=n-8_r_q3_n-2}\hfill\qr{1b}{j=potencias\&m=9\&c=n-2_e3_s_n5}\hfill
\qr{1c}{j=potencias\&m=9\&c=r_n3_e2_s_q2_a_n-1_s_n4_c}\hfill\qr{2d}{j=potencias\&m=8\&c=n-2_e2_p-3_s_n1}\hfill
\qr{reto}{j=potencias\&t=1\&c=r_n2_e2}\hfill\qr{reto}{j=potencias\&t=1\&c=n-2_e2}}
\end{dorso}
\end{document}
